\documentclass{beamer}
\usepackage[utf8]{inputenc}
\usepackage[T1]{fontenc}
\usepackage{graphicx}
\usepackage{amsmath}
\usepackage{amssymb}
\usepackage{mathtools}
\usepackage{array}
\usepackage{qtree}
\usepackage{hyperref}
\usepackage[slovene]{babel}
\usepackage{color}
\definecolor{red}{rgb}{1,0.4,0.4}
\definecolor{green}{rgb}{0.4,0.75,0.4}
\definecolor{blue}{rgb}{0.4,0.4,0.75}
\definecolor{yellow}{rgb}{0.75,0.75,0.4}
\definecolor{grey}{rgb}{0.8,0.8,0.8}
<<<<<<< HEAD
\setbeamertemplate{navigation symbols}{}
\usetheme{metropolis}
\theoremstyle{plain}
\newtheorem{izrek}{Izrek}
\theoremstyle{plain}
\newtheorem{opomba}{Opomba}
\let\oldblock\block
\let\endoldblock\endblock
=======

\setbeamertemplate{navigation symbols}{}
\usetheme{metropolis}

\theoremstyle{plain}
\newtheorem{izrek}{Izrek}

\theoremstyle{plain}
\newtheorem{opomba}{Opomba}

\let\oldblock\block
\let\endoldblock\endblock

>>>>>>> 4bef1e4a4ce04f95cfb275cb27d304305e4641d9
\renewenvironment{block}[1]
    {\begin{oldblock}{#1}
        \smallskip
    }
    { 
    \end{oldblock}
    }
<<<<<<< HEAD
\let\oldalertblock\alertblock
\let\endoldalertblock\endalertblock
=======


\let\oldalertblock\alertblock
\let\endoldalertblock\endalertblock

>>>>>>> 4bef1e4a4ce04f95cfb275cb27d304305e4641d9
\renewenvironment{alertblock}[1]
    {\begin{oldalertblock}{#1}
        \smallskip
    }
    { 
    \end{oldalertblock}
    }
<<<<<<< HEAD
\let\oldexampleblock\exampleblock
\let\endoldexampleblock\endexampleblock
=======


\let\oldexampleblock\exampleblock
\let\endoldexampleblock\endexampleblock


>>>>>>> 4bef1e4a4ce04f95cfb275cb27d304305e4641d9
\renewenvironment{exampleblock}[1]
    {\begin{oldexampleblock}{#1}
        \smallskip
    }
    { 
    \end{oldexampleblock}
    }
<<<<<<< HEAD
=======


>>>>>>> 4bef1e4a4ce04f95cfb275cb27d304305e4641d9
% Naloga 1.3.1: Za dokument uporabite razred `beamer'.
% Ne dodajajte nastavitve za velikost pisave, kot je bila v datoteki `5-prosojnice.tex`.

% Naloga 1.3.2: vključite paket `predavanja'.
<<<<<<< HEAD
=======

>>>>>>> 4bef1e4a4ce04f95cfb275cb27d304305e4641d9
% Naloga 1.3.3: definirajte okolje `izrek'.
% Namig: z iskanjem po datotekah (Ctrl+Shift+F oz. Cmd+Shift+F) 
% poiščite niz `{izrek}'.
\title{Matematični izrazi in uporaba paketa \texttt{beamer}}
\subtitle{Vaje pri predmetu \emph{Računalniški praktikum}}
\institute{Fakulteta za matematiko in fiziko}
\date{}
\usepackage{tikz}
\usetikzlibrary{intersections}
\begin{document}

\frame{\titlepage}

\begin{frame}
\frametitle{Kratek pregled}    
\tableofcontents
\end{frame}

% Naloga 1.3.4: pripravite naslovno stran z vsebino:
% - naslov: Matematični izrazi in uporaba paketa \texttt{beamer}
% - podnaslov: Vaje pri predmetu \emph{Računalniški praktikum}
% - ustanova: Fakulteta za matematiko in fiziko
% - datum: naj se ne izpiše; to dosežete z ukazom \date{}.
% Zgornje podatke nastavite z ukazi kot v dokumentih razreda `article`.
% Več o tem, kako se naredi naslovno stran, si preberite na naslovu:
% https://www.overleaf.com/learn/latex/Beamer
% To stran preberite do vključno razdelka "Creating a table of contents".
% Ukaz `\titlepage` deluje podobno kot ukaz `\maketitle`, ki ste ga že srečali.
% Naloga 1.3.5: pripravite kazalo vsebine.
% 1. Naslov prosojnice s kazalom vsebine naj bo "Kratek pregled".
% 2. Z neobveznim argumentom `pausesections' (v oglatih oklepajih)
%    določite, da se kazalo vsebine odkriva postopoma.
%    Poglejte, kako deluje ta ukaz.
% 3. Ker ni videti preveč lepo, neobvezni argument zakomentirajte.
<<<<<<< HEAD
<<<<<<< HEAD

\section{Paket \texttt{beamer}}
% Naslov prosojnice lahko napišemo tudi kot argument okolja `frame`.
\begin{frame}{Posebnosti prosojnic}
	% Naloga 2.3.1:
	% Dodajte ukaze, ki bodo poskrbeli, da se bo prosojnica odkrivala postopoma,
	% tako kot v datoteki prosojnice-resitev.pdf

	 Za prosojnice je značilna uporaba okolja \texttt{frame},
    s katerim definiramo posamezno prosojnico,
    \pause
    postopno odkrivanje prosojnic,
    \pause
    ter nekateri drugi ukazi, ki jih najdemo v paketu \texttt{beamer}.
    \pause
    \begin{exampleblock}{Primer}
        Verjetno ste že opazili, da za naslovno prosojnico niste uporabili
        ukaza \texttt{maketitle}, ampak ukaz \texttt{titlepage}.
    \end{exampleblock}
\end{frame}

\begin{frame}{Poudarjeni bloki}
	% Naloga 2.3.2:
	% Oblikujte poudarjena bloka z opombo in opozorilom.

		\begin{block}{Opomba}
        Okolja za poudarjene bloke so \texttt{block}, \texttt{exampleblock} in \texttt{alertblock}.
    \end{block}
    
    \pause
    
    \begin{alertblock}{Pozor!}
        Začetek poudarjenega bloka (ukaz \texttt{begin}) vedno sprejme
        dva argumenta: okolje in naslov bloka.
        Drugi argument (za naslov) je lahko prazen.
    \end{alertblock} 

\end{frame}

\begin{frame}{Tudi v predstavitvah lahko pišemo izreke in dokaze}
	% Naloga 2.3.3:
	% Oblikujte okolje itemize, tako da se bo njegova vsebina postopoma odkrivala.
	% Ne smete uporabiti ukaza `pause'.
	% Beseda `največje' naj bo poudarjena šele na četrti podprosojnici.

	\begin{izrek}
	   Praštevil je neskončno mnogo.
	\end{izrek}
	\begin{proof}
	   Denimo, da je praštevil končno mnogo.
		% Z neobveznim argumentom [<+->] okolja itemize določimo,
		% da se elementi naštevanja odkrivajo postopoma.
		   \begin{itemize}[<+->]
            \item Naj bo \(p\) \alert<4>{največje} praštevilo.
            \item Naj bo \(q\) produkt števil \(1\), \(2\), \ldots, \(p\).
            \item Število \(q+1\) ni deljivo z nobenim praštevilom, torej je \(q+1\) praštevilo.
            \item To je protislovje, saj je \(q+1>p\). \qedhere
	   \end{itemize}
	\end{proof}
 \end{frame}
 

=======
\section{Paket \texttt{beamer}}
% Naslov prosojnice lahko napišemo tudi kot argument okolja `frame`.
\begin{frame}{Posebnosti prosojnic}
	% Naloga 2.3.1:
	% Dodajte ukaze, ki bodo poskrbeli, da se bo prosojnica odkrivala postopoma,
	% tako kot v datoteki prosojnice-resitev.pdf

	 Za prosojnice je značilna uporaba okolja \texttt{frame},
    s katerim definiramo posamezno prosojnico,
    \pause
    postopno odkrivanje prosojnic,
    \pause
    ter nekateri drugi ukazi, ki jih najdemo v paketu \texttt{beamer}.
    \pause
    \begin{exampleblock}{Primer}
        Verjetno ste že opazili, da za naslovno prosojnico niste uporabili
        ukaza \texttt{maketitle}, ampak ukaz \texttt{titlepage}.
    \end{exampleblock}
\end{frame}

\begin{frame}{Poudarjeni bloki}
	% Naloga 2.3.2:
	% Oblikujte poudarjena bloka z opombo in opozorilom.

		\begin{block}{Opomba}
        Okolja za poudarjene bloke so \texttt{block}, \texttt{exampleblock} in \texttt{alertblock}.
    \end{block}
    
    \pause
    
    \begin{alertblock}{Pozor!}
        Začetek poudarjenega bloka (ukaz \texttt{begin}) vedno sprejme
        dva argumenta: okolje in naslov bloka.
        Drugi argument (za naslov) je lahko prazen.
    \end{alertblock} 

\end{frame}

\begin{frame}{Tudi v predstavitvah lahko pišemo izreke in dokaze}
	% Naloga 2.3.3:
	% Oblikujte okolje itemize, tako da se bo njegova vsebina postopoma odkrivala.
	% Ne smete uporabiti ukaza `pause'.
	% Beseda `največje' naj bo poudarjena šele na četrti podprosojnici.

	\begin{izrek}
	   Praštevil je neskončno mnogo.
	\end{izrek}
	\begin{proof}
	   Denimo, da je praštevil končno mnogo.
		% Z neobveznim argumentom [<+->] okolja itemize določimo,
		% da se elementi naštevanja odkrivajo postopoma.
		   \begin{itemize}[<+->]
            \item Naj bo \(p\) \alert<4>{največje} praštevilo.
            \item Naj bo \(q\) produkt števil \(1\), \(2\), \ldots, \(p\).
            \item Število \(q+1\) ni deljivo z nobenim praštevilom, torej je \(q+1\) praštevilo.
            \item To je protislovje, saj je \(q+1>p\). \qedhere
	   \end{itemize}
	\end{proof}
 \end{frame}
 

>>>>>>> e11faaa93e7f9f9d49c88c797eb1154b4dc920c5
=======

\section{Paket \texttt{beamer}}
% Naslov prosojnice lahko napišemo tudi kot argument okolja `frame`.
\begin{frame}{Posebnosti prosojnic}
	% Naloga 2.3.1:
	% Dodajte ukaze, ki bodo poskrbeli, da se bo prosojnica odkrivala postopoma,
	% tako kot v datoteki prosojnice-resitev.pdf

	 Za prosojnice je značilna uporaba okolja \texttt{frame},
    s katerim definiramo posamezno prosojnico,
    \pause
    postopno odkrivanje prosojnic,
    \pause
    ter nekateri drugi ukazi, ki jih najdemo v paketu \texttt{beamer}.
    \pause
    \begin{exampleblock}{Primer}
        Verjetno ste že opazili, da za naslovno prosojnico niste uporabili
        ukaza \texttt{maketitle}, ampak ukaz \texttt{titlepage}.
    \end{exampleblock}
\end{frame}

\begin{frame}{Poudarjeni bloki}
	% Naloga 2.3.2:
	% Oblikujte poudarjena bloka z opombo in opozorilom.

		\begin{block}{Opomba}
        Okolja za poudarjene bloke so \texttt{block}, \texttt{exampleblock} in \texttt{alertblock}.
    \end{block}
    
    \pause
    
    \begin{alertblock}{Pozor!}
        Začetek poudarjenega bloka (ukaz \texttt{begin}) vedno sprejme
        dva argumenta: okolje in naslov bloka.
        Drugi argument (za naslov) je lahko prazen.
    \end{alertblock} 

\end{frame}

\begin{frame}{Tudi v predstavitvah lahko pišemo izreke in dokaze}
	% Naloga 2.3.3:
	% Oblikujte okolje itemize, tako da se bo njegova vsebina postopoma odkrivala.
	% Ne smete uporabiti ukaza `pause'.
	% Beseda `največje' naj bo poudarjena šele na četrti podprosojnici.

	\begin{izrek}
	   Praštevil je neskončno mnogo.
	\end{izrek}
	\begin{proof}
	   Denimo, da je praštevil končno mnogo.
		% Z neobveznim argumentom [<+->] okolja itemize določimo,
		% da se elementi naštevanja odkrivajo postopoma.
		   \begin{itemize}[<+->]
            \item Naj bo \(p\) \alert<4>{največje} praštevilo.
            \item Naj bo \(q\) produkt števil \(1\), \(2\), \ldots, \(p\).
            \item Število \(q+1\) ni deljivo z nobenim praštevilom, torej je \(q+1\) praštevilo.
            \item To je protislovje, saj je \(q+1>p\). \qedhere
	   \end{itemize}
	\end{proof}
 \end{frame}
 

>>>>>>> 4bef1e4a4ce04f95cfb275cb27d304305e4641d9

\section{Paketa \texttt{amsmath} in \texttt{amsfonts}}
\begin{frame}{Postopno odkrivanje v okolju align}
	% Naloga 3.2:
	% Za pravilno postopno odkrivanje morate na enem mestu uporabiti ukaz `only',
	% na dveh pa ukaz `onslide'.
	% Namig: prva vrstica enačbe se od druge podprosojnice naprej nadaljuje
	% z vsebino druge vrstice. To dosežete tako, da z ukazom `only' samo na prvi
	% podprosojnici postavite del prve vrstice od znaka & do znakov \\,
	% vključno z njimi:
	%     (a+b)^n \only<1>{&= \ldots \\}
	% Na poznejših podprosojnicah tako ni ne treh pik ne konca vrstice,
	% zato se za (a+b)^n izpiše naslednja vrstica.
	Dokaži \emph{binomsko formulo}: za vsaki realni števili \(a\) in \(b\) in za vsako naravno število
\begin{align*}
(a+b)^n \only<1>{&= \ldots \\}
\onslide<2->{&= (a+b) (a+b) \dots (a+b) \\}
\onslide<3->{&= a^n + n a^{n-1} b + \dots + \binom{n}{k} a^{n-k} b^k + \dots + n a b^{n-1} + b^n \\}
\onslide<3->{&= \sum_{k=0}^{n} \binom{n}{k} a^{n-k} b^k}
\end{align*}
\end{frame}


\section{Stolpci in slike}
\begin{frame}{Konstrukcija pravokotnice}
	\begin{columns}
	\begin{column}{0.55\textwidth}
	\begin{itemize}[<+->]
		\item Dana sta premica \(p\) in točka \(T\).
		\item Nariši lok \(k\) s središčem v \(T\).
		\item Lok \(k\) seče premico \(p\) v točkah \(A\) in \(B\).
		\item Nariši lok \(m\) s središčem v \(A\).
		\item Nariši lok \(n\) s središčem v \(B\) in z enakim polmerom.
		\item Loka se sečeta v točki \(C\).
		\item Premica \(TC\) je pravokotna na \(p\).
	\end{itemize}
\end{column}
\begin{column}{0.45\textwidth}
	\centering
	\includegraphics<1>[width=50mm]{slike/fig-1.png}
	\includegraphics<2>[width=50mm]{slike/fig-2.png}
	\includegraphics<3>[width=50mm]{slike/fig-3.png}
	\includegraphics<4>[width=50mm]{slike/fig-4.png}
	\includegraphics<5>[width=50mm]{slike/fig-5.png}
	\includegraphics<6>[width=50mm]{slike/fig-6.png}
	\includegraphics<7>[width=50mm]{slike/fig-7.png}
\end{column}
\end{columns}
\end{frame}


<<<<<<< HEAD
<<<<<<< HEAD
\section{Paket \texttt{beamer} in tabele}
\begin{frame}{Odkrivanje tabele po vrsticah}
    Včasih pride prav, da tabelo odkrivamo postopoma po vrsticah.
    \begin{center}
        \begin{tabular}{c|cccc}
            Oznaka & A & B & C & D \\ \hline
            \pause
            X & 1 & 2 & 3 & 4 \\
            \pause
            Y & 3 & 4 & 5 & 6 \\
            \pause
            Z & 5 & 6 & 7 & 8
        \end{tabular}
    \end{center}
\end{frame}

\begin{frame}{Odkrivanje tabele po stolpcih}
    Tabelo lahko odkrivamo tudi po stolpcih, čeprav ni najlažje.
    \begin{center}
        \begin{tabular}{c|cccc}
            Oznaka & \onslide<2->A & \onslide<2->B & \onslide<3->C & \onslide<3->D \\ \hline
            X & \onslide<2->1 & \onslide<2->2 & \onslide<3->3 & \onslide<3->4 \\
            Y & \onslide<2->3 & \onslide<2->4 & \onslide<3->5 & \onslide<3->6 \\
            Z & \onslide<2->5 & \onslide<2->6 & \onslide<3->7 & \onslide<3->8
        \end{tabular}
    \end{center}
\end{frame}
\section{Risanje s paketom TikZ}
\begin{frame}{Konstrukcija pravokotnice}
	\centering
	\begin{tikzpicture}
		\coordinate (T) at (0, 1.4);

		\draw<1->[blue, very thick, name path=p] (-3, 0) -- (3, 0) node[right] {$p$};
		\fill<1->[blue] (T) circle (2pt) node[above left] {$T$};

		\onslide<2->{
			\draw[gray, name path=k] (T) ++(210:2) arc (210:330:2);
			\path[name intersections={of=p and k, by={A, B}}];
			\draw (A) circle (1.5pt) node[below left] {$A$};
			\draw (B) circle (1.5pt) node[below right] {$B$};
		}

		\onslide<3->{
			\draw[gray, name path=m] (A) ++(-70:2) arc (-70:70:2);
			\draw[gray, name path=n] (B) ++(250:2) arc (250:110:2);
			\path[name intersections={of=m and n, by={C, D}}];
		}

		\onslide<4->{
			\draw (C) circle (1.5pt) node[below right] {$C$};
		}

		\onslide<5->{
			\draw[red, very thick] (T) -- (C);
		}
	\end{tikzpicture}
\end{frame}
\end{document}

=======
\section{Paket beamer in tabele}
=======
\section{Paket \texttt{beamer} in tabele}
>>>>>>> 4bef1e4a4ce04f95cfb275cb27d304305e4641d9
\begin{frame}{Odkrivanje tabele po vrsticah}
    Včasih pride prav, da tabelo odkrivamo postopoma po vrsticah.
    \begin{center}
        \begin{tabular}{c|cccc}
            Oznaka & A & B & C & D \\ \hline
            \pause
            X & 1 & 2 & 3 & 4 \\
            \pause
            Y & 3 & 4 & 5 & 6 \\
            \pause
            Z & 5 & 6 & 7 & 8
        \end{tabular}
    \end{center}
\end{frame}

\begin{frame}{Odkrivanje tabele po stolpcih}
    Tabelo lahko odkrivamo tudi po stolpcih, čeprav ni najlažje.
    \begin{center}
        \begin{tabular}{c|cccc}
            Oznaka & \onslide<2->A & \onslide<2->B & \onslide<3->C & \onslide<3->D \\ \hline
            X & \onslide<2->1 & \onslide<2->2 & \onslide<3->3 & \onslide<3->4 \\
            Y & \onslide<2->3 & \onslide<2->4 & \onslide<3->5 & \onslide<3->6 \\
            Z & \onslide<2->5 & \onslide<2->6 & \onslide<3->7 & \onslide<3->8
        \end{tabular}
    \end{center}
\end{frame}
\section{Risanje s paketom TikZ}
\begin{frame}{Konstrukcija pravokotnice}
	\centering
	\begin{tikzpicture}
		\coordinate (T) at (0, 1.4);

		\draw<1->[blue, very thick, name path=p] (-3, 0) -- (3, 0) node[right] {$p$};
		\fill<1->[blue] (T) circle (2pt) node[above left] {$T$};

		\onslide<2->{
			\draw[gray, name path=k] (T) ++(210:2) arc (210:330:2);
			\path[name intersections={of=p and k, by={A, B}}];
			\draw (A) circle (1.5pt) node[below left] {$A$};
			\draw (B) circle (1.5pt) node[below right] {$B$};
		}

		\onslide<3->{
			\draw[gray, name path=m] (A) ++(-70:2) arc (-70:70:2);
			\draw[gray, name path=n] (B) ++(250:2) arc (250:110:2);
			\path[name intersections={of=m and n, by={C, D}}];
		}

		\onslide<4->{
			\draw (C) circle (1.5pt) node[below right] {$C$};
		}

		\onslide<5->{
			\draw[red, very thick] (T) -- (C);
		}
	\end{tikzpicture}
\end{frame}
<<<<<<< HEAD

\end{document}
>>>>>>> e11faaa93e7f9f9d49c88c797eb1154b4dc920c5
=======
\end{document}
>>>>>>> 4bef1e4a4ce04f95cfb275cb27d304305e4641d9
